\section{嘉靖隆庆的朝廷、法律与文化连结}

本节以朝廷程序、法律语言、家庭与性别、城市文化和海外联系为问题线索。正史、实录、律例与奏议通常更容易保存官员、男性精英和国家分类；女性、家庭成员、工匠、城市居民与海外中介则常仅在墓志、契约、地方志、碑刻、航海记录或案件中留下片段。材料本身的偏向，应当成为叙事的一部分。

\subsection{礼制裁断与法律语言}

\eventanchor{MM-1531-ONE-WHIP-PROPOSAL}{明·1531（嘉靖十年）}
\paragraph{1531（嘉靖十年）：桂萼等讨论合并赋役、折纳银两的方案，显示一条鞭式改革的早期政策脉络} 此事把权力、法律或文化关系推进到可辨认的节点。账本的证据限度是奏议可证当时的改革主张；它不是全国统一法令，亦不能把后来的万历制度直接倒投为嘉靖十年的既成事实。可资核查的材料为《世宗实录》嘉靖十年户部条；《明史·食货志》；桂萼奏议及地方赋役材料。它们能够说明制度如何表述自身、某种命令或交涉如何被记录，却未必能直接复原家庭内部关系、性别分工、城市日常或海外行动者的意图。事件发生的条件涉及里甲役法、实物征解与土地兼并使各地赋役负担和征收成本日益失衡，官员寻求折银和合并名目的办法；其后留下的制度与社会后果为相关方案在不同地区以不同速度试行或搁置；至万历时的清丈、一条鞭推行不能视为其线性完成。这也要求将正式名义、实际执行和当事人的处境分别讨论。

\eventanchor{ML-1532-001}{明·1532（嘉靖十一年）}
\paragraph{1532（嘉靖十一年）：嘉靖倭寇大袭：广州巡抚因铲除倭寇不力被罢免} 此事把权力、法律或文化关系推进到可辨认的节点。账本的证据限度是编年候选的年序可由官修与后出材料交叉；具体月日、参与者、战果或数字必须回查所列材料，不能将单一叙事当作定论。可资核查的材料为《世宗实录》嘉靖十一年条；《明史·兵志》及相关列传；边镇、地方志或对方材料。它们能够说明制度如何表述自身、某种命令或交涉如何被记录，却未必能直接复原家庭内部关系、性别分工、城市日常或海外行动者的意图。事件发生的条件涉及前期边防、地方武装或跨区域资源调动的压力推动了该行动；其后留下的制度与社会后果为它改变了后续调兵、补给或地方秩序的条件；战果和伤亡须分开核对。译写候选以编年材料定位；实录记载反映朝廷选择，崇祯期须另以奏疏、地方及敌对政权材料交叉。

\eventanchor{ML-1532-002}{明·1532（嘉靖十一年）}
\paragraph{1532（嘉靖十一年）：看到了一颗非常不吉利的彗星} 此事把权力、法律或文化关系推进到可辨认的节点。账本的证据限度是编年候选的年序可由官修与后出材料交叉；具体月日、参与者、战果或数字必须回查所列材料，不能将单一叙事当作定论。可资核查的材料为《世宗实录》嘉靖十一年条；《明史·本纪》及相关列传；诏令、奏疏或会典版本。它们能够说明制度如何表述自身、某种命令或交涉如何被记录，却未必能直接复原家庭内部关系、性别分工、城市日常或海外行动者的意图。事件发生的条件涉及继承、官僚协调或皇权运作的压力使问题进入朝廷程序；其后留下的制度与社会后果为它影响后续人事与政策议程；动机和派系标签不能由单一叙事裁定。译写候选以编年材料定位；实录记载反映朝廷选择，崇祯期须另以奏疏、地方及敌对政权材料交叉。

\eventanchor{ML-1532-003}{明·1532（嘉靖十一年）}
\paragraph{1532（嘉靖十一年）：北边警报、边镇防务和西南处置的本年军政材料标记边疆压力，敌我称谓与军数应回到原文。（材料核查重点：诏令或奏议的形成与发受链）} 此事把权力、法律或文化关系推进到可辨认的节点。账本的证据限度是来源导向的年度桥接条目：本年材料可定位相关政务线索；具体月日、发文机关、地域和执行结果仍须逐项回查，不能以制度文本或奏报替代实际效果。可资核查的材料为《世宗实录》嘉靖十一年条；《明史·兵志》及相关列传；边镇、地方志或对方材料。它们能够说明制度如何表述自身、某种命令或交涉如何被记录，却未必能直接复原家庭内部关系、性别分工、城市日常或海外行动者的意图。事件发生的条件涉及前期边防、地方武装或跨区域资源调动的压力推动了该行动；其后留下的制度与社会后果为它改变了后续调兵、补给或地方秩序的条件；战果和伤亡须分开核对。桥接条目只标示可回查的年度问题与材料链；实录有中央选择性，崇祯期尤其需要奏疏、地方、朝鲜及满文材料交叉。；本条特别核对诏令或奏议的形成与发受链。

\eventanchor{ML-1532-004}{明·1532（嘉靖十一年）}
\paragraph{1532（嘉靖十一年）：海禁、朝贡和沿海港口管理的本年材料显示官方海上政策，禁令范围与实际航行须分别调查} 此事把权力、法律或文化关系推进到可辨认的节点。账本的证据限度是来源导向的年度桥接条目：本年材料可定位相关政务线索；具体月日、发文机关、地域和执行结果仍须逐项回查，不能以制度文本或奏报替代实际效果。可资核查的材料为《世宗实录》嘉靖十一年条；《明史·外国传》；《朝鲜王朝实录》、琉球《历代宝案》或海外档案。它们能够说明制度如何表述自身、某种命令或交涉如何被记录，却未必能直接复原家庭内部关系、性别分工、城市日常或海外行动者的意图。事件发生的条件涉及朝贡名义、边境安全与港口贸易的利益使交涉持续发生；其后留下的制度与社会后果为它为后续册封、互市或海防安排留下文书与跨政权比较线索。桥接条目只标示可回查的年度问题与材料链；实录有中央选择性，崇祯期尤其需要奏疏、地方、朝鲜及满文材料交叉。

\eventanchor{ML-1532-005}{明·1532（嘉靖十一年）}
\paragraph{1532（嘉靖十一年）：北边警报、边镇防务和西南处置的本年军政材料标记边疆压力，敌我称谓与军数应回到原文。（材料核查重点：报称信息的来源、抄传与行政分类）} 此事把权力、法律或文化关系推进到可辨认的节点。账本的证据限度是来源导向的年度桥接条目：本年材料可定位相关政务线索；具体月日、发文机关、地域和执行结果仍须逐项回查，不能以制度文本或奏报替代实际效果。可资核查的材料为《世宗实录》嘉靖十一年条；《明史·兵志》及相关列传；边镇、地方志或对方材料。它们能够说明制度如何表述自身、某种命令或交涉如何被记录，却未必能直接复原家庭内部关系、性别分工、城市日常或海外行动者的意图。事件发生的条件涉及前期边防、地方武装或跨区域资源调动的压力推动了该行动；其后留下的制度与社会后果为它改变了后续调兵、补给或地方秩序的条件；战果和伤亡须分开核对。桥接条目只标示可回查的年度问题与材料链；实录有中央选择性，崇祯期尤其需要奏疏、地方、朝鲜及满文材料交叉。；本条特别核对报称信息的来源、抄传与行政分类。

\subsection{家庭、宫廷与性别材料}

\eventanchor{ML-1533-001}{明·1533（嘉靖十二年）十月24日}
\paragraph{1533（嘉靖十二年）十月24日：大同卫军叛乱，被镇压} 此事把权力、法律或文化关系推进到可辨认的节点。账本的证据限度是编年候选的年序可由官修与后出材料交叉；具体月日、参与者、战果或数字必须回查所列材料，不能将单一叙事当作定论。可资核查的材料为《世宗实录》嘉靖十二年条；《明史·本纪》及相关列传；诏令、奏疏或会典版本。它们能够说明制度如何表述自身、某种命令或交涉如何被记录，却未必能直接复原家庭内部关系、性别分工、城市日常或海外行动者的意图。事件发生的条件涉及继承、官僚协调或皇权运作的压力使问题进入朝廷程序；其后留下的制度与社会后果为它影响后续人事与政策议程；动机和派系标签不能由单一叙事裁定。译写候选以编年材料定位；实录记载反映朝廷选择，崇祯期须另以奏疏、地方及敌对政权材料交叉。

\eventanchor{ML-1533-002}{明·1533（嘉靖十二年）}
\paragraph{1533（嘉靖十二年）：海禁、朝贡和沿海港口管理的本年材料显示官方海上政策，禁令范围与实际航行须分别调查} 此事把权力、法律或文化关系推进到可辨认的节点。账本的证据限度是来源导向的年度桥接条目：本年材料可定位相关政务线索；具体月日、发文机关、地域和执行结果仍须逐项回查，不能以制度文本或奏报替代实际效果。可资核查的材料为《世宗实录》嘉靖十二年条；《明史·外国传》；《朝鲜王朝实录》、琉球《历代宝案》或海外档案。它们能够说明制度如何表述自身、某种命令或交涉如何被记录，却未必能直接复原家庭内部关系、性别分工、城市日常或海外行动者的意图。事件发生的条件涉及朝贡名义、边境安全与港口贸易的利益使交涉持续发生；其后留下的制度与社会后果为它为后续册封、互市或海防安排留下文书与跨政权比较线索。桥接条目只标示可回查的年度问题与材料链；实录有中央选择性，崇祯期尤其需要奏疏、地方、朝鲜及满文材料交叉。

\eventanchor{ML-1533-003}{明·1533（嘉靖十二年）}
\paragraph{1533（嘉靖十二年）：赋役折色、盐课、河工或田土核查的本年文书为财政改革史提供节点，制度名称不可跨区套用。（材料核查重点：诏令或奏议的形成与发受链）} 此事把权力、法律或文化关系推进到可辨认的节点。账本的证据限度是来源导向的年度桥接条目：本年材料可定位相关政务线索；具体月日、发文机关、地域和执行结果仍须逐项回查，不能以制度文本或奏报替代实际效果。可资核查的材料为《世宗实录》嘉靖十二年条；《明会典》相关条目；地方志、黄册/契约/河工材料。它们能够说明制度如何表述自身、某种命令或交涉如何被记录，却未必能直接复原家庭内部关系、性别分工、城市日常或海外行动者的意图。事件发生的条件涉及财政征解、交通工程或货币支付的压力使相关措施进入政务；其后留下的制度与社会后果为其法定安排不等于地方执行，后续须以地域材料追踪负担与成效。桥接条目只标示可回查的年度问题与材料链；实录有中央选择性，崇祯期尤其需要奏疏、地方、朝鲜及满文材料交叉。；本条特别核对诏令或奏议的形成与发受链。

\eventanchor{ML-1533-004}{明·1533（嘉靖十二年）}
\paragraph{1533（嘉靖十二年）：祀典、学校、出版或讲学的本年材料记录文化制度活动，规范话语不等于社会普遍认同} 此事把权力、法律或文化关系推进到可辨认的节点。账本的证据限度是来源导向的年度桥接条目：本年材料可定位相关政务线索；具体月日、发文机关、地域和执行结果仍须逐项回查，不能以制度文本或奏报替代实际效果。可资核查的材料为《世宗实录》嘉靖十二年条；《明史·选举志》或《艺文志》；文集、碑刻或书目材料。它们能够说明制度如何表述自身、某种命令或交涉如何被记录，却未必能直接复原家庭内部关系、性别分工、城市日常或海外行动者的意图。事件发生的条件涉及制度、教育或知识传播的需要推动了相关行动或文本形成；其后留下的制度与社会后果为它提供制度和传播史的锚点，不能据此推断社会整体接受。桥接条目只标示可回查的年度问题与材料链；实录有中央选择性，崇祯期尤其需要奏疏、地方、朝鲜及满文材料交叉。

\eventanchor{ML-1533-005}{明·1533（嘉靖十二年）}
\paragraph{1533（嘉靖十二年）：赋役折色、盐课、河工或田土核查的本年文书为财政改革史提供节点，制度名称不可跨区套用。（材料核查重点：报称信息的来源、抄传与行政分类）} 此事把权力、法律或文化关系推进到可辨认的节点。账本的证据限度是来源导向的年度桥接条目：本年材料可定位相关政务线索；具体月日、发文机关、地域和执行结果仍须逐项回查，不能以制度文本或奏报替代实际效果。可资核查的材料为《世宗实录》嘉靖十二年条；《明会典》相关条目；地方志、黄册/契约/河工材料。它们能够说明制度如何表述自身、某种命令或交涉如何被记录，却未必能直接复原家庭内部关系、性别分工、城市日常或海外行动者的意图。事件发生的条件涉及财政征解、交通工程或货币支付的压力使相关措施进入政务；其后留下的制度与社会后果为其法定安排不等于地方执行，后续须以地域材料追踪负担与成效。桥接条目只标示可回查的年度问题与材料链；实录有中央选择性，崇祯期尤其需要奏疏、地方、朝鲜及满文材料交叉。；本条特别核对报称信息的来源、抄传与行政分类。

\eventanchor{ML-1534-001}{明·1534（嘉靖十三年）}
\paragraph{1534（嘉靖十三年）：嘉靖倭寇突袭：俘获一名率领50多艘大船的海盗} 此事把权力、法律或文化关系推进到可辨认的节点。账本的证据限度是编年候选的年序可由官修与后出材料交叉；具体月日、参与者、战果或数字必须回查所列材料，不能将单一叙事当作定论。可资核查的材料为《世宗实录》嘉靖十三年条；《明史·兵志》及相关列传；边镇、地方志或对方材料。它们能够说明制度如何表述自身、某种命令或交涉如何被记录，却未必能直接复原家庭内部关系、性别分工、城市日常或海外行动者的意图。事件发生的条件涉及前期边防、地方武装或跨区域资源调动的压力推动了该行动；其后留下的制度与社会后果为它改变了后续调兵、补给或地方秩序的条件；战果和伤亡须分开核对。译写候选以编年材料定位；实录记载反映朝廷选择，崇祯期须另以奏疏、地方及敌对政权材料交叉。

\subsection{都市、宗教与文化空间}

\eventanchor{ML-1534-002}{明·1534（嘉靖十三年）}
\paragraph{1534（嘉靖十三年）：嘉靖皇帝停止例行朝见} 此事把权力、法律或文化关系推进到可辨认的节点。账本的证据限度是编年候选的年序可由官修与后出材料交叉；具体月日、参与者、战果或数字必须回查所列材料，不能将单一叙事当作定论。可资核查的材料为《世宗实录》嘉靖十三年条；《明史·本纪》及相关列传；诏令、奏疏或会典版本。它们能够说明制度如何表述自身、某种命令或交涉如何被记录，却未必能直接复原家庭内部关系、性别分工、城市日常或海外行动者的意图。事件发生的条件涉及继承、官僚协调或皇权运作的压力使问题进入朝廷程序；其后留下的制度与社会后果为它影响后续人事与政策议程；动机和派系标签不能由单一叙事裁定。译写候选以编年材料定位；实录记载反映朝廷选择，崇祯期须另以奏疏、地方及敌对政权材料交叉。

\eventanchor{ML-1534-003}{明·1534（嘉靖十三年）}
\paragraph{1534（嘉靖十三年）：海禁、朝贡和沿海港口管理的本年材料显示官方海上政策，禁令范围与实际航行须分别调查。（材料核查重点：诏令或奏议的形成与发受链）} 此事把权力、法律或文化关系推进到可辨认的节点。账本的证据限度是来源导向的年度桥接条目：本年材料可定位相关政务线索；具体月日、发文机关、地域和执行结果仍须逐项回查，不能以制度文本或奏报替代实际效果。可资核查的材料为《世宗实录》嘉靖十三年条；《明史·外国传》；《朝鲜王朝实录》、琉球《历代宝案》或海外档案。它们能够说明制度如何表述自身、某种命令或交涉如何被记录，却未必能直接复原家庭内部关系、性别分工、城市日常或海外行动者的意图。事件发生的条件涉及朝贡名义、边境安全与港口贸易的利益使交涉持续发生；其后留下的制度与社会后果为它为后续册封、互市或海防安排留下文书与跨政权比较线索。桥接条目只标示可回查的年度问题与材料链；实录有中央选择性，崇祯期尤其需要奏疏、地方、朝鲜及满文材料交叉。；本条特别核对诏令或奏议的形成与发受链。

\eventanchor{ML-1534-004}{明·1534（嘉靖十三年）}
\paragraph{1534（嘉靖十三年）：灾异、赈济及地方工程的本年报告可与州县材料互校，救济命令不自动证明救济到达} 此事把权力、法律或文化关系推进到可辨认的节点。账本的证据限度是来源导向的年度桥接条目：本年材料可定位相关政务线索；具体月日、发文机关、地域和执行结果仍须逐项回查，不能以制度文本或奏报替代实际效果。可资核查的材料为《世宗实录》嘉靖十三年条；《明史·五行志》；受灾府县地方志与奏疏。它们能够说明制度如何表述自身、某种命令或交涉如何被记录，却未必能直接复原家庭内部关系、性别分工、城市日常或海外行动者的意图。事件发生的条件涉及自然风险与地方报告促成朝廷或地方的应对记录；其后留下的制度与社会后果为赈济、迁徙和损失的实际范围仍须以受灾地区材料分别核验。桥接条目只标示可回查的年度问题与材料链；实录有中央选择性，崇祯期尤其需要奏疏、地方、朝鲜及满文材料交叉。

\eventanchor{ML-1534-005}{明·1534（嘉靖十三年）}
\paragraph{1534（嘉靖十三年）：海禁、朝贡和沿海港口管理的本年材料显示官方海上政策，禁令范围与实际航行须分别调查。（材料核查重点：报称信息的来源、抄传与行政分类）} 此事把权力、法律或文化关系推进到可辨认的节点。账本的证据限度是来源导向的年度桥接条目：本年材料可定位相关政务线索；具体月日、发文机关、地域和执行结果仍须逐项回查，不能以制度文本或奏报替代实际效果。可资核查的材料为《世宗实录》嘉靖十三年条；《明史·外国传》；《朝鲜王朝实录》、琉球《历代宝案》或海外档案。它们能够说明制度如何表述自身、某种命令或交涉如何被记录，却未必能直接复原家庭内部关系、性别分工、城市日常或海外行动者的意图。事件发生的条件涉及朝贡名义、边境安全与港口贸易的利益使交涉持续发生；其后留下的制度与社会后果为它为后续册封、互市或海防安排留下文书与跨政权比较线索。桥接条目只标示可回查的年度问题与材料链；实录有中央选择性，崇祯期尤其需要奏疏、地方、朝鲜及满文材料交叉。；本条特别核对报称信息的来源、抄传与行政分类。

\eventanchor{ML-1535-001}{明·1535（嘉靖十四年）}
\paragraph{1535（嘉靖十四年）：辽东、广宁守军叛乱，被镇压} 此事把权力、法律或文化关系推进到可辨认的节点。账本的证据限度是编年候选的年序可由官修与后出材料交叉；具体月日、参与者、战果或数字必须回查所列材料，不能将单一叙事当作定论。可资核查的材料为《世宗实录》嘉靖十四年条；《明史·本纪》及相关列传；诏令、奏疏或会典版本。它们能够说明制度如何表述自身、某种命令或交涉如何被记录，却未必能直接复原家庭内部关系、性别分工、城市日常或海外行动者的意图。事件发生的条件涉及继承、官僚协调或皇权运作的压力使问题进入朝廷程序；其后留下的制度与社会后果为它影响后续人事与政策议程；动机和派系标签不能由单一叙事裁定。译写候选以编年材料定位；实录记载反映朝廷选择，崇祯期须另以奏疏、地方及敌对政权材料交叉。

\eventanchor{ML-1535-002}{明·1535（嘉靖十四年）}
\paragraph{1535（嘉靖十四年）：祀典、学校、出版或讲学的本年材料记录文化制度活动，规范话语不等于社会普遍认同。（材料核查重点：诏令或奏议的形成与发受链）} 此事把权力、法律或文化关系推进到可辨认的节点。账本的证据限度是来源导向的年度桥接条目：本年材料可定位相关政务线索；具体月日、发文机关、地域和执行结果仍须逐项回查，不能以制度文本或奏报替代实际效果。可资核查的材料为《世宗实录》嘉靖十四年条；《明史·选举志》或《艺文志》；文集、碑刻或书目材料。它们能够说明制度如何表述自身、某种命令或交涉如何被记录，却未必能直接复原家庭内部关系、性别分工、城市日常或海外行动者的意图。事件发生的条件涉及制度、教育或知识传播的需要推动了相关行动或文本形成；其后留下的制度与社会后果为它提供制度和传播史的锚点，不能据此推断社会整体接受。桥接条目只标示可回查的年度问题与材料链；实录有中央选择性，崇祯期尤其需要奏疏、地方、朝鲜及满文材料交叉。；本条特别核对诏令或奏议的形成与发受链。

\subsection{港口、朝贡与海外联系}

\eventanchor{ML-1535-003}{明·1535（嘉靖十四年）}
\paragraph{1535（嘉靖十四年）：祀典、学校、出版或讲学的本年材料记录文化制度活动，规范话语不等于社会普遍认同。（材料核查重点：报称信息的来源、抄传与行政分类）} 此事把权力、法律或文化关系推进到可辨认的节点。账本的证据限度是来源导向的年度桥接条目：本年材料可定位相关政务线索；具体月日、发文机关、地域和执行结果仍须逐项回查，不能以制度文本或奏报替代实际效果。可资核查的材料为《世宗实录》嘉靖十四年条；《明史·选举志》或《艺文志》；文集、碑刻或书目材料。它们能够说明制度如何表述自身、某种命令或交涉如何被记录，却未必能直接复原家庭内部关系、性别分工、城市日常或海外行动者的意图。事件发生的条件涉及制度、教育或知识传播的需要推动了相关行动或文本形成；其后留下的制度与社会后果为它提供制度和传播史的锚点，不能据此推断社会整体接受。桥接条目只标示可回查的年度问题与材料链；实录有中央选择性，崇祯期尤其需要奏疏、地方、朝鲜及满文材料交叉。；本条特别核对报称信息的来源、抄传与行政分类。

\eventanchor{ML-1535-004}{明·1535（嘉靖十四年）}
\paragraph{1535（嘉靖十四年）：嘉靖朝礼制、内阁与人事的本年诏令和奏议形成中枢协调线索，留中与施行之间应予区分} 此事把权力、法律或文化关系推进到可辨认的节点。账本的证据限度是来源导向的年度桥接条目：本年材料可定位相关政务线索；具体月日、发文机关、地域和执行结果仍须逐项回查，不能以制度文本或奏报替代实际效果。可资核查的材料为《世宗实录》嘉靖十四年条；《明史·本纪》及相关列传；诏令、奏疏或会典版本。它们能够说明制度如何表述自身、某种命令或交涉如何被记录，却未必能直接复原家庭内部关系、性别分工、城市日常或海外行动者的意图。事件发生的条件涉及继承、官僚协调或皇权运作的压力使问题进入朝廷程序；其后留下的制度与社会后果为它影响后续人事与政策议程；动机和派系标签不能由单一叙事裁定。桥接条目只标示可回查的年度问题与材料链；实录有中央选择性，崇祯期尤其需要奏疏、地方、朝鲜及满文材料交叉。

\eventanchor{ML-1535-005}{明·1535（嘉靖十四年）}
\paragraph{1535（嘉靖十四年）：祀典、学校、出版或讲学的本年材料记录文化制度活动，规范话语不等于社会普遍认同。（材料核查重点：同年地方材料所见的执行范围）} 此事把权力、法律或文化关系推进到可辨认的节点。账本的证据限度是来源导向的年度桥接条目：本年材料可定位相关政务线索；具体月日、发文机关、地域和执行结果仍须逐项回查，不能以制度文本或奏报替代实际效果。可资核查的材料为《世宗实录》嘉靖十四年条；《明史·选举志》或《艺文志》；文集、碑刻或书目材料。它们能够说明制度如何表述自身、某种命令或交涉如何被记录，却未必能直接复原家庭内部关系、性别分工、城市日常或海外行动者的意图。事件发生的条件涉及制度、教育或知识传播的需要推动了相关行动或文本形成；其后留下的制度与社会后果为它提供制度和传播史的锚点，不能据此推断社会整体接受。桥接条目只标示可回查的年度问题与材料链；实录有中央选择性，崇祯期尤其需要奏疏、地方、朝鲜及满文材料交叉。；本条特别核对同年地方材料所见的执行范围。

\eventanchor{ML-1536-001}{明·1536（嘉靖十五年）}
\paragraph{1536（嘉靖十五年）：蒙古人袭击山西但被击退} 此事把权力、法律或文化关系推进到可辨认的节点。账本的证据限度是编年候选的年序可由官修与后出材料交叉；具体月日、参与者、战果或数字必须回查所列材料，不能将单一叙事当作定论。可资核查的材料为《世宗实录》嘉靖十五年条；《明史·兵志》及相关列传；边镇、地方志或对方材料。它们能够说明制度如何表述自身、某种命令或交涉如何被记录，却未必能直接复原家庭内部关系、性别分工、城市日常或海外行动者的意图。事件发生的条件涉及前期边防、地方武装或跨区域资源调动的压力推动了该行动；其后留下的制度与社会后果为它改变了后续调兵、补给或地方秩序的条件；战果和伤亡须分开核对。译写候选以编年材料定位；实录记载反映朝廷选择，崇祯期须另以奏疏、地方及敌对政权材料交叉。

\eventanchor{ML-1536-002}{明·1536（嘉靖十五年）}
\paragraph{1536（嘉靖十五年）：灾异、赈济及地方工程的本年报告可与州县材料互校，救济命令不自动证明救济到达。（材料核查重点：诏令或奏议的形成与发受链）} 此事把权力、法律或文化关系推进到可辨认的节点。账本的证据限度是来源导向的年度桥接条目：本年材料可定位相关政务线索；具体月日、发文机关、地域和执行结果仍须逐项回查，不能以制度文本或奏报替代实际效果。可资核查的材料为《世宗实录》嘉靖十五年条；《明史·五行志》；受灾府县地方志与奏疏。它们能够说明制度如何表述自身、某种命令或交涉如何被记录，却未必能直接复原家庭内部关系、性别分工、城市日常或海外行动者的意图。事件发生的条件涉及自然风险与地方报告促成朝廷或地方的应对记录；其后留下的制度与社会后果为赈济、迁徙和损失的实际范围仍须以受灾地区材料分别核验。桥接条目只标示可回查的年度问题与材料链；实录有中央选择性，崇祯期尤其需要奏疏、地方、朝鲜及满文材料交叉。；本条特别核对诏令或奏议的形成与发受链。

\eventanchor{ML-1536-003}{明·1536（嘉靖十五年）}
\paragraph{1536（嘉靖十五年）：灾异、赈济及地方工程的本年报告可与州县材料互校，救济命令不自动证明救济到达。（材料核查重点：报称信息的来源、抄传与行政分类）} 此事把权力、法律或文化关系推进到可辨认的节点。账本的证据限度是来源导向的年度桥接条目：本年材料可定位相关政务线索；具体月日、发文机关、地域和执行结果仍须逐项回查，不能以制度文本或奏报替代实际效果。可资核查的材料为《世宗实录》嘉靖十五年条；《明史·五行志》；受灾府县地方志与奏疏。它们能够说明制度如何表述自身、某种命令或交涉如何被记录，却未必能直接复原家庭内部关系、性别分工、城市日常或海外行动者的意图。事件发生的条件涉及自然风险与地方报告促成朝廷或地方的应对记录；其后留下的制度与社会后果为赈济、迁徙和损失的实际范围仍须以受灾地区材料分别核验。桥接条目只标示可回查的年度问题与材料链；实录有中央选择性，崇祯期尤其需要奏疏、地方、朝鲜及满文材料交叉。；本条特别核对报称信息的来源、抄传与行政分类。
